\begin{frame}{Comparing Sequence Backbones}
{\fontsize{10}{12}\selectfont We compare how four network designs use time-ordered data.\par}
\vspace{0pt}
\begingroup
\fontsize{9.2}{10}\selectfont
\setlength{\tabcolsep}{5pt}
\renewcommand{\arraystretch}{0.96}
\begin{tabular}{@{}>{\raggedright\arraybackslash}p{.31\linewidth}>{\raggedright\arraybackslash}p{.25\linewidth}>{\raggedright\arraybackslash}p{\dimexpr.44\linewidth-20pt\relax}@{}}
\textbf{Main idea} & \textbf{Advantage} & \textbf{Weakness} \\ [0pt]\hline
\textbf{Vanilla recurrent neural network} \RefLink{bengio1994}{1}\newline
The network reads time-ordered data.
& It carries information from earlier steps.
& The gradient vanishes in long sequences, so the network forgets early steps. \\[2pt]\hline
\textbf{Long short-term memory} \RefLink{hochreiter1997}{2}\newline
Gates control what the memory keeps and forgets.
& The gates help the network retain earlier information.
& It processes steps one by one, so training is slow. It also has more parameters. \\[2pt]\hline
\textbf{Convolutional / temporal convolutional network} \RefLink{bai2018}{3}\newline
Filters read a fixed window of nearby steps.
& The filters process steps in parallel.
& The network sees only as far as its effective window. \\[2pt]\hline
\textbf{Transformer} \RefLink{vaswani2017}{4}\newline
Attention links any two time steps directly.
& It connects distant steps without a recurrent chain.
& Its attention cost grows with the square of sequence length. It needs a time encoding. \\[2pt]\hline
\end{tabular}
\endgroup
\end{frame}
